% ══════════════════════════════════════════════════════
% PART 3: RAILWAY DEPLOYMENT (VERY DETAILED)
% ══════════════════════════════════════════════════════

\chapter{Phase 3 --- Deploy on Railway.app}

\section{What is Railway?}

\textbf{Railway.app} (\url{https://railway.app}) is a cloud platform that runs your code 24/7. Think of it as a computer in the cloud that never turns off.

\begin{table}[H]
\centering
\caption{Railway at a Glance}
\begin{tabularx}{\textwidth}{|l|X|}
\hline
\textbf{Feature} & \textbf{Details} \\
\hline
Website & \url{https://railway.app} \\
\hline
What it does & Runs your Python bot 24/7 in the cloud \\
\hline
How it works & Reads your code from GitHub and runs it \\
\hline
Free trial & \$5 free credit (enough for 2--3 weeks of testing) \\
\hline
After trial & \$5/month (Hobby plan) + usage (typically \$2--5/month) \\
\hline
Payment & Credit card required after trial \\
\hline
Difficulty & Very easy --- no Linux/server knowledge needed \\
\hline
\end{tabularx}
\end{table}

\begin{infobox}
\textbf{Railway vs VPS:} Railway is like hiring someone to manage your computer for you. You just give them the code, and they handle everything (installing Python, running the bot, restarting if it crashes). A VPS (Chapter 4) is like buying your own computer --- you manage everything yourself.
\end{infobox}

\section{Step 1: Create a Railway Account}

\begin{stepbox}{Create Your Account}
\begin{enumerate}[leftmargin=1.5cm]
    \item Open your browser and go to \url{https://railway.app}
    \item You will see a landing page with a big button
    \item Click \textbf{``Login''} (top right) or \textbf{``Start a New Project''}
    \item You will see login options:
    \begin{itemize}
        \item \textbf{Login with GitHub} $\leftarrow$ \textbf{Choose this one!}
        \item Login with Email
    \end{itemize}
    \item Click \textbf{``Login with GitHub''}
    \item A GitHub authorization page appears
    \item Click \textbf{``Authorize Railway''}
    \item You are now logged into Railway!
\end{enumerate}
\end{stepbox}

\begin{tipbox}
We use ``Login with GitHub'' because Railway needs access to your GitHub repository to read the bot code. This is the easiest way to connect them.
\end{tipbox}

\section{Step 2: Create a New Project}

\begin{stepbox}{Start a New Project}
\begin{enumerate}[leftmargin=1.5cm]
    \item After logging in, you'll see the Railway dashboard
    \item Click \textbf{``+ New Project''} (it's a prominent button)
    \item A menu appears with options:
    \begin{itemize}
        \item Deploy from GitHub repo $\leftarrow$ \textbf{Choose this one!}
        \item Deploy a Template
        \item Provision a Database
        \item Empty Project
    \end{itemize}
    \item Click \textbf{``Deploy from GitHub repo''}
\end{enumerate}
\end{stepbox}

\section{Step 3: Connect Your GitHub Repository}

\begin{stepbox}{Select Your Repository}
\begin{enumerate}[leftmargin=1.5cm]
    \item Railway shows a list of your GitHub repositories
    \item If you see \texttt{active-dream-bot} in the list:
    \begin{itemize}
        \item Click on it to select it
    \end{itemize}
    \item If you do NOT see it (because it's private):
    \begin{enumerate}
        \item Click \textbf{``Configure GitHub App''} at the bottom
        \item A GitHub page opens asking which repositories Railway can access
        \item Under ``Repository access'', select:
        \begin{itemize}
            \item \textbf{``Only select repositories''} $\leftarrow$ Choose this for security
        \end{itemize}
        \item Click the dropdown and select \texttt{active-dream-bot}
        \item Click \textbf{``Save''}
        \item Go back to Railway and you should now see the repository
        \item Click on \texttt{active-dream-bot}
    \end{enumerate}
\end{enumerate}
\end{stepbox}

\begin{warningbox}
\textbf{If Railway tries to deploy immediately}, it will \textbf{fail} because we haven't added the environment variables yet. That's okay! We'll fix it in the next step.
\end{warningbox}

\section{Step 4: Add Environment Variables}

This is the \textbf{most important step}. You need to add all the secret values from your \texttt{.env} file into Railway's settings.

\begin{stepbox}{Adding Variables}
\begin{enumerate}[leftmargin=1.5cm]
    \item After selecting your repository, Railway creates a ``service'' for your bot
    \item You should see your project dashboard with a box/card representing your service
    \item \textbf{Click on the service card} (it might be called ``active-dream-bot'')
    \item You'll see a panel on the right side with tabs
    \item Click the \textbf{``Variables''} tab
    \item You'll see an empty variables section
    \item Now you need to add each variable \textbf{one by one}:
\end{enumerate}
\end{stepbox}

For each variable, click \textbf{``+ New Variable''} (or the ``Add'' button) and enter the name and value:

\begin{table}[H]
\centering
\caption{Variables to Add in Railway}
\begin{tabularx}{\textwidth}{|l|X|}
\hline
\textbf{Variable Name} & \textbf{Value (use client's actual values)} \\
\hline
\texttt{BOT\_TOKEN} & \texttt{7123456789:AAFx8Kz3mN9bR2wQ...} (from BotFather) \\
\hline
\texttt{ADMIN\_IDS} & \texttt{987654321} (client's Telegram User ID) \\
\hline
\texttt{FIVESIM\_API\_KEY} & \texttt{eyJhbGciOiJSUzUxMi...} (from 5sim.net) \\
\hline
\texttt{OTP\_GROUP\_ID} & \texttt{-1001234567890} (from @RawDataBot) \\
\hline
\texttt{OTP\_GROUP\_LINK} & \texttt{https://t.me/ActiveDreamOTP} \\
\hline
\texttt{CHANNEL\_LINK} & \texttt{https://t.me/ActiveDreamChannel} \\
\hline
\texttt{BOT\_NAME} & \texttt{Active Dream} \\
\hline
\texttt{OTP\_CHECK\_INTERVAL} & \texttt{5} \\
\hline
\texttt{OTP\_MAX\_WAIT} & \texttt{300} \\
\hline
\texttt{PYTHONIOENCODING} & \texttt{utf-8} \\
\hline
\end{tabularx}
\end{table}

\begin{importantbox}
\textbf{Add \texttt{PYTHONIOENCODING=utf-8} as a variable!} Without this, the bot may crash on Railway due to emoji encoding issues. This is a common issue that is easy to miss.
\end{importantbox}

\begin{tipbox}
\textbf{Bulk Add Shortcut:} Some versions of Railway allow you to paste multiple variables at once in the format:
\begin{lstlisting}[style=termstyle]
BOT_TOKEN=7123456789:AAFx8Kz3mN9bR2wQ...
ADMIN_IDS=987654321
FIVESIM_API_KEY=eyJhbGci...
OTP_GROUP_ID=-1001234567890
\end{lstlisting}
Try clicking ``RAW Editor'' if available --- you can paste all variables at once.
\end{tipbox}

\section{Step 5: Configure the Start Command}

Railway needs to know \textbf{how to start your bot}. If you created the \texttt{Procfile} in Phase 2, Railway should detect it automatically. But let's make sure:

\begin{stepbox}{Verify/Set Start Command}
\begin{enumerate}[leftmargin=1.5cm]
    \item Click on your service in Railway
    \item Go to the \textbf{``Settings''} tab
    \item Look for \textbf{``Start Command''} or \textbf{``Custom Start Command''}
    \item If it's empty, type:
    \begin{lstlisting}[style=termstyle]
python bot.py
    \end{lstlisting}
    \item Make sure the \textbf{``Build Command''} or install command is:
    \begin{lstlisting}[style=termstyle]
pip install -r requirements.txt
    \end{lstlisting}
    (Railway usually detects this automatically for Python projects)
\end{enumerate}
\end{stepbox}

\section{Step 6: Deploy}

\begin{stepbox}{Trigger Deployment}
\begin{enumerate}[leftmargin=1.5cm]
    \item After adding all variables, Railway should \textbf{automatically redeploy}
    \item If it doesn't, click \textbf{``Deploy''} or \textbf{``Redeploy''} button
    \item Railway will now:
    \begin{enumerate}
        \item Read your code from GitHub
        \item Install Python
        \item Install dependencies from \texttt{requirements.txt}
        \item Run \texttt{python bot.py}
    \end{enumerate}
    \item You can watch the progress in the \textbf{``Deployments''} tab
    \item The deployment takes about \textbf{1--3 minutes}
\end{enumerate}
\end{stepbox}

\section{Step 7: Check the Logs}

After deployment, you need to verify the bot started successfully:

\begin{stepbox}{Viewing Logs}
\begin{enumerate}[leftmargin=1.5cm]
    \item Click on your service in Railway
    \item Go to the \textbf{``Deployments''} tab
    \item Click on the \textbf{latest deployment} (should show a green dot if successful)
    \item Click \textbf{``View Logs''} or the log icon
    \item You should see output like:
    \begin{lstlisting}[style=termstyle]
[BOT] Active Dream is starting...
INFO - Bot initialized. Database ready.
INFO - Application started
    \end{lstlisting}
    \item If you see these messages, \textbf{the bot is live!}
\end{enumerate}
\end{stepbox}

\subsection{Common Errors in Logs}

\begin{table}[H]
\centering
\caption{Troubleshooting Railway Deployment Errors}
\begin{tabularx}{\textwidth}{|l|X|}
\hline
\textbf{Error Message} & \textbf{Solution} \\
\hline
\texttt{BOT\_TOKEN is not set} & You forgot to add the \texttt{BOT\_TOKEN} variable. Go to Variables tab and add it. \\
\hline
\texttt{ADMIN\_IDS is not set} & Add the \texttt{ADMIN\_IDS} variable with the client's User ID. \\
\hline
\texttt{ModuleNotFoundError} & The \texttt{requirements.txt} file is missing or incorrect. Make sure it's in the root folder. \\
\hline
\texttt{UnicodeEncodeError} & Add \texttt{PYTHONIOENCODING=utf-8} as a variable. \\
\hline
\texttt{401 Unauthorized} & The \texttt{BOT\_TOKEN} is invalid. Double-check with the client. \\
\hline
\texttt{Build failed} & Check if \texttt{requirements.txt} exists and has correct package names. \\
\hline
\end{tabularx}
\end{table}

\section{Step 8: Test the Bot}

\begin{stepbox}{Verification Steps}
\begin{enumerate}[leftmargin=1.5cm]
    \item Open Telegram on your phone
    \item Search for the bot's username (e.g., \texttt{@ActiveDream\_OTP\_Bot})
    \item Tap \textbf{Start} or send \texttt{/start}
    \item You should see the \textbf{``Request Access''} button
    \item Click it --- you should see ``Access Pending''
    \item The \textbf{client} (admin) should receive a notification with Approve/Reject
    \item After approval, send \texttt{/start} again --- full menu should appear
    \item Try \textbf{Get Number} $\rightarrow$ select any service $\rightarrow$ select any country
    \item A number should appear, and OTP should arrive within seconds (mock) or minutes (real)
\end{enumerate}
\end{stepbox}

\begin{successbox}
\textbf{If the bot responds on Telegram, congratulations! Your bot is deployed and running 24/7 on Railway!}

The bot will continue running even when you close your browser or turn off your computer. Railway keeps it running in the cloud.
\end{successbox}

\section{Railway Dashboard Overview}

After deployment, your Railway dashboard is your control panel:

\begin{table}[H]
\centering
\caption{Railway Dashboard Tabs}
\begin{tabularx}{\textwidth}{|l|X|}
\hline
\textbf{Tab} & \textbf{What It Shows} \\
\hline
\textbf{Deployments} & History of all deployments. Click any to see logs. \\
\hline
\textbf{Variables} & Your environment variables (.env values). Edit here. \\
\hline
\textbf{Settings} & Start command, region, build settings. \\
\hline
\textbf{Metrics} & CPU and memory usage graphs. \\
\hline
\textbf{Networking} & Not needed for Telegram bots (they don't need a URL). \\
\hline
\end{tabularx}
\end{table}

\section{How to Update the Bot Later}

If you make changes to the code in the future:

\begin{enumerate}[leftmargin=1.5cm]
    \item Edit the files on your PC
    \item Open terminal and run:
    \begin{lstlisting}[style=termstyle]
cd "D:\Dart Project\active_dream_bot"
git add .
git commit -m "Description of changes"
git push
    \end{lstlisting}
    \item Railway will \textbf{automatically detect the new code} and redeploy
    \item The bot restarts with the new code within 1--2 minutes
\end{enumerate}

\begin{tipbox}
This is one of the biggest advantages of Railway: every time you push code to GitHub, Railway automatically updates the running bot. No need to SSH into a server or manually restart anything.
\end{tipbox}

\section{Railway Pricing After Free Trial}

\begin{table}[H]
\centering
\caption{Railway Cost Breakdown}
\begin{tabularx}{\textwidth}{|l|c|X|}
\hline
\textbf{Item} & \textbf{Cost} & \textbf{Notes} \\
\hline
Hobby Plan & \$5/month & Base subscription fee \\
\hline
Compute usage & \$1--3/month & Based on CPU/memory; a Telegram bot uses very little \\
\hline
\textbf{Total} & \textbf{\$6--8/month} & Approximately 700--950 BDT \\
\hline
\end{tabularx}
\end{table}

\begin{warningbox}
\textbf{Important Note About Railway and SQLite:} Railway's file system is \textbf{ephemeral} --- meaning if Railway redeploys your app, the database file (\texttt{active\_dream.db}) gets \textbf{deleted and recreated}. This means:
\begin{itemize}
    \item All user data, balances, and transaction history are lost on every redeploy
    \item This is fine for testing and short-term use
    \item For long-term production use, you should either:
    \begin{enumerate}
        \item Use a VPS instead (Chapter 4) --- database persists permanently
        \item Add a Railway PostgreSQL database (more complex, but data persists)
    \end{enumerate}
\end{itemize}
\textbf{Recommendation:} Use Railway for initial deployment and client demonstration. If the client uses the bot long-term with many users, migrate to a VPS (Chapter 4).
\end{warningbox}
